\section{Introduction}
\label{sec:introduction}

A perturbation of delayed network coupling can vanish from the equation for
the stationary network average at every fixed history and still change the
local passage through a fold.  The perturbation changes transverse histories,
which re-enter the average equation through heterogeneous node nonlinearities.
We study this effect in a fast--slow network and quantify its first variation
uniformly in the number of nodes.

Let \(P_N\) be row-stochastic with stationary probability \(\pi_N\).
The network has the form
\begin{align}
 \dot v(t)={}&f_N(v(t),w(t))+D(P_N-\Id)v(t)\notag\\
 &+\delta^2K\sum_{k=0}^m B_{k,N}(\eta)
   \left\{v(t)-v\left(t-\frac{\theta_k}{\delta}\right)\right\},
                                                        \label{eq:intro-network-form}\\
 \dot w(t)={}&\delta^2\{\pi_N^Tv(t)-a\},              \label{eq:intro-recovery-form}
\end{align}
where \(0\le\theta_0<\cdots<\theta_m\), \(a=1+\delta^2\nu\), and
\(0<\delta\ll1\).  The polynomial field \(f_N\), specified in
\cref{sec:model-results}, has positive quadratic coefficients \(c_{i,N}\)
with fixed stationary average \(\alpha=\pi_N^Tc_N>0\).
We write \(B_{k,N}(\eta)=B_{k,N}+\Delta B_{k,N}\) and impose
\begin{equation}
 \pi_N^T\Delta B_{k,N}=0\quad(0\le k\le m),
 \qquad \sum_{k=0}^m\Delta B_{k,N}=0.                \label{eq:intro-neutrality}
\end{equation}
Thus the sum of the delayed-coupling matrices is unchanged, and each
perturbation is annihilated by the stationary projection.  If
\(\mathcal F_{N,\delta,\nu,\eta}\) denotes the retarded functional
differential equation (RFDE) vector field and
\(\Pi_N(v,w)=(\pi_N^Tv,w)\), then
\begin{equation}
 \Pi_N\mathcal F_{N,\delta,\nu,\eta}(\Phi)
 =\Pi_N\mathcal F_{N,\delta,\nu,0}(\Phi)              \label{eq:intro-projected-identity}
\end{equation}
for every full history \(\Phi\).  This identity compares vector fields at the
same history; it does not close the dynamics of the average.  The question is
how the change of histories affects a scalar matching condition near the fold.

\subsection{Local fold matching in relation to earlier theory}

Finite-dimensional canard theory describes passage between attracting and
repelling slow manifolds near a fold
\citep{krupa2001relaxation,krupa2001extending}.  For delay equations,
\citet{krupa2016explosion} combine slow-manifold persistence, local center
manifolds, and fast connections to obtain canard explosions, with an
application to the delayed van der Pol system.  Their companion study
\citep{krupa2016complex} analyzes the more complex oscillations generated as
the delay increases.  The general local theory of
\citet[Theorem~4.5 and Corollary~4.6]{avitabile2020local} constructs
parameter-dependent center manifolds under spectral-gap and regularity
hypotheses for infinite-dimensional fast systems coupled to finitely many
slow variables.  The reduced equations retain the slow--fast structure,
allowing finite-dimensional analysis near the loss of normal hyperbolicity.

More recently, \citet{zhang2026high} study a delayed van der Pol equation
whose feedback gain is of the same order as the time-scale separation
parameter.  They reduce the equation to a planar flow on a nonlocal center
manifold and use a nonlinear time transformation to obtain high-order
approximations of the canard critical value and manifold.  Thus a reduction
that retains delay-dependent history information is already part of the
local theory.  Our question concerns the derivative with respect to the
constrained matrix directions in \eqref{eq:intro-neutrality}.  We construct
the compatible histories and the finite-interval matching function for this
network, and bound that derivative uniformly over a growing family of RFDE
phase spaces.  In physical time the delays are \(\theta_k/\delta\); in fold
time \(s=\delta t\) they are the fixed translations \(\theta_k\).

The principal result is \cref{thm:rw-local-fold-matching}.  Under the uniform
coefficient bounds and the Dobrushin mixing condition stated in
\cref{sec:model-results}, there is a two-dimensional locally invariant
manifold of compatible RFDE histories near the fold.  Two reduced
boundary-value problems define a scalar matching function
\(D_N^{\rm fin}(\delta,\nu,\eta)\).  Its expansion and first derivatives
satisfy
\begin{align}
 \frac{D_N^{\rm fin}}{\delta}
 &=\sqrt{2\pi}(\nu-\nu_{0,N})+O(\delta),\notag\\
 \partial_\nu D_N^{\rm fin}
 &=\delta\sqrt{2\pi}+O(\delta^2),\notag\\
 \bigl\|D_\eta D_N^{\rm fin}
       +\delta^2\sqrt{2\pi}\Lambda_N\bigr\|
 &\le C\{\delta^3+\delta^2\|\eta\|\},              \label{eq:intro-local-matching}
\end{align}
where \(\nu_{0,N}\) is the explicit fold center in \cref{eq:rw-fold-center}.
Consequently there is a unique nearby matching zero
\(\nu_N^{\rm fin}(\delta,\eta)\), and the parameter
\(\mu_N^{\rm fin}=\delta^2\nu_N^{\rm fin}\) satisfies
\begin{equation}
 \bigl\|D_\eta\mu_N^{\rm fin}-\delta^3\Lambda_N\bigr\|
 \le C\{\delta^4+\delta^3\|\eta\|\}.                \label{eq:intro-local-root-derivative}
\end{equation}
The neighborhoods, constants, and remainders are independent of \(N\), in
the norms of \cref{sec:model-results}.  This zero belongs to the specified
finite boundary-value problems.  Their trajectories represent solutions of
the original RFDE while they and their delayed histories remain in the fold
neighborhood.  The auxiliary completion fixes the finite matching problem;
its zero is not identified with a global connection.

\subsection{The delay moment and the proof mechanism}

The coefficient in \eqref{eq:intro-local-matching} factors through the
network-transverse space \(E_N=\ker\pi_N^T\).  Set
\[
 P_{\perp,N}=\Id-\mathbf1\pi_N^T,
 \qquad A_N=D(P_N-\Id)|_{E_N}.
\]
The first delay moment produces the transverse forcing
\begin{equation}
 \mathsf S_N(\eta)
 =\frac{K}{2\alpha}P_{\perp,N}
   \left(\sum_{k=0}^m\theta_k\Delta B_{k,N}\right)\mathbf1,
                                                        \label{eq:intro-delay-moment}
\end{equation}
and the sensitivity functional is
\begin{equation}
 \Lambda_N=\mathfrak r_NA_N^{-1}\mathsf S_N,
 \qquad
 \mathfrak r_N(h)=-\frac1\alpha\pi_N^T\diag(c_N)h.
                                                        \label{eq:intro-response-covector}
\end{equation}
Here \(A_N^{-1}\mathsf S_N\) is the leading perturbation derivative of the
transverse graph coefficient.  Heterogeneous curvature returns this displacement to the
collective equation, and pairing with its one-dimensional Gaussian cokernel
gives \(\mathfrak r_N\).  Direct cancellation in
\eqref{eq:intro-projected-identity} therefore need not cancel the derivative
of the matching function.

Two estimates make this calculation quantitative.  First, two distinct delay
locations allow every vector in \(E_N\) to be realized by opposite rank-one
matrix perturbations at those locations.  The right inverse of
\(\mathsf S_N\) is norm-optimal in the specified norms
(\cref{thm:rw-first-moment-map}); with only one delay location after equal
delays are merged, \(\mathsf S_N\) vanishes.  For a prescribed matrix
pattern \(\mathcal U_N\), the attainable space for opposite two-layer
directions is \(\{R\mathbf1:R\in\mathcal U_N\}\), up to a nonzero scalar.
Preserving each layer's row sums makes this first-moment response vanish;
generator-supported patterns can retain the full range
(\cref{thm:rw-structured-delay-range}).  Second, the Dobrushin condition
\begin{equation}
 \tau(P_N)\le1-\gamma,\qquad\gamma>0,                \label{eq:intro-dobrushin-gap}
\end{equation}
gives \(\|A_N^{-1}\|\le(D\gamma)^{-1}\) on \(E_N\), equipped with the
oscillation norm.  This standard contraction estimate
\citep{gaubert2015dobrushin} yields the uniform bound for \(\Lambda_N\) in
\cref{cor:rw-fold-fredholm-coefficient}.

The main analytic step is to establish \eqref{eq:intro-local-matching} for
actual compatible histories.  A variation-of-constants construction uses the
transverse contraction to obtain the invariant graph and its mixed state
and parameter derivatives on the growing fold interval
\(S_\delta=\sqrt{2p_0\log(1/\delta)}\), with \(p_0>4\).
Reconstructing each history along the reduced flow retains the fixed delay
translations.  Differentiation of this graph gives
\(A_N^{-1}\mathsf S_N\); the Gaussian adjoint then evaluates its collective
contribution.  Uniform estimates for the two finite boundary-value problems
control the moving traces, graph remainders, and omitted Gaussian tails,
passing the coefficient to the first derivatives of \(D_N^{\rm fin}\).
The rank-one construction and the Fredholm elimination support this argument;
they are elementary and standard steps, respectively.

For a concrete example, take
\[
 \pi=\tfrac13(1,1,1)^T,\qquad P=\mathbf1\pi^T,\qquad
 c=(1-\sigma,1,1+\sigma)^T,\qquad s=(-1,0,1)^T,
\]
where \(0<\sigma<1\).  Choose \((\theta_0,\theta_1)=(1,2)\) and
\(R_0=s\pi^T\), \(R_1=-s\pi^T\).  Then
\(\pi^TR_k=0\), \(R_0+R_1=0\), but
\(\theta_0R_0\mathbf1+\theta_1R_1\mathbf1=-s\ne0\).
For \(D=1\), \(K=2\), and \(\alpha=1\), one has \(A=-\Id\), so
\begin{equation}
 \Lambda(R_0,R_1)=-\pi^T\diag(c)s=-\frac{2\sigma}{3}.
                                                        \label{eq:intro-three-node-response}
\end{equation}
If \(B_0=B_1=P/2\), the perturbed layers \(B_k+\zeta R_k\) remain
nonnegative with the same full support for \(|\zeta|<1/2\), and their sum
remains \(P\).  The distinct quadratic coefficients exclude nontrivial
coordinate-equality invariant subspaces.  In
\cref{prop:rw-growing-families}, a uniformly mixing family realizes a
nonzero coefficient for every \(N\) with uniformly bounded perturbation
directions, while a lazy cycle shows why uniformity fails for general local
mixing.  The computations in \cref{sec:fold-passage} illustrate these
coefficients using finite-section diagnostics.

A heteroclinic interpretation requires additional information about the
vector field away from the fold.  In \cref{sec:heteroclinic-connection} we
specify a modified recovery equation and assume a scalar defining function
for the relevant unstable/stable manifold intersection, together with a
uniform \(C^1\) comparison to \(D_N^{\rm fin}\).
Under these hypotheses, \cref{cor:rw-conditional-connection} transfers
\(\Lambda_N\) to the derivative of a connection parameter.  The global
comparison is an assumption; general RFDE connection theory and Lin-type
reductions \citep{hale1986heteroclinic,malletparet1999fredholm,hupkes2009lin}
provide a setting in which it can be studied.

Section~2 states the model and the local results.  Section~3 proves the
first-moment range and Fredholm formulas, and Section~4 constructs the
invariant histories and finite-interval matching function.  Section~5 gives
the conditional connection application.  Section~6 records the scope of the
local result; Appendix~A contains the detailed graph and matching estimates.
